%% ma: L10-B17-0001
%% lop: 10
%% chuong: 6
%% bai: 17
%% dang: 10.17.1
%% loai: TN4
%% muc: TH
%% dap_an: "C"
%% nguon: "Ảnh giáo viên gửi 24/10/2026 (M06): Kỳ thi thử đợt IV, Thi thử TN THPT 2027 (Biên soạn: Aimbet), Phần 1 Câu 10"
%% kien_thuc: "Bài 17 (dấu của tam thức bậc hai, xét dấu tích)"
%% trang_thai: chinh-thuc
\begin{ex}
Hàm số $f(x)=\left(x^2-3x+2\right)\left(x^2-4\right)$ nhận giá trị âm khi và chỉ khi
\choice{$x\in(-2;1)\cup(2;3)$}{$x\in(1;2)$}{\True $x\in(-2;1)$}{$x\in(-\infty;-2)\cup(1;2)$}
\loigiai{$f(x)=(x-1)(x-2)^2(x+2)$. Với $x\ne2$ thì $(x-2)^2>0$, nên $f(x)<0\Leftrightarrow(x-1)(x+2)<0\Leftrightarrow-2<x<1$.}
\end{ex}
